\section{From even weights to full profinite characters}\label{sec:characters}
This section explains why rational even-weight indecomposables produce
characters of the full gentle group which vanish on the pentagon group.
The rational argument does not identify a mod-$p$ pentagon kernel with a
reduction of its characteristic-zero kernel.

We use the Deligne--Drinfeld theorem in the precise form proved by
OpenAI and streamlined by Willwacher \cite[Theorem 1.1 and the discussion
of (7), pp.~1--2]{Willwacher}: $\grt(\Q)$ is generated, as a graded Lie
algebra, by one element in each odd weight at least three.  Brown's
injectivity theorem supplies freeness, although generation alone is enough
here.  As $\grt\subset W$, it follows that
\begin{equation}
 \ker P_m=\grt(\Q)_m\subset[W,W]_m
       \qquad(m\text{ even}).\label{eq:even-pentagon}
\end{equation}
Indeed an even-weight Lie word in odd-weight generators has at least two
letters and is a bracket of two positive-weight Lie words.  The same
inclusion holds after scalar extension to $\Q_p$.

\begin{theorem}[Even-weight characters]\label{thm:even-characters}
Let $p\geq5$ be prime, let $m\geq2$ be even, and put
$h_m=\dim_\Q(W/[W,W])_m$.  For every $k\geq1$ such that
$(p-1)p^{k-1}$ divides $m$, there is a continuous surjection
\[
 \GTz/\overline{\langle\!\langle\GT\rangle\!\rangle}
       \longrightarrow(\Z/p^k\Z)^{h_m}.
\]
The map is already onto on the global multiplier-one subgroup of $\GTz$.
In particular weight sixteen gives a quotient $(\Z/17\Z)^{38}$.
\end{theorem}

The proof needs a full realization theorem.  We spell out its local
part, and use Theorem~\ref{thm:peripheral-lifting} for the passage to
the full profinite group.

\begin{lemma}[Realization of special directions]\label{lem:realization}
For $p\geq5$, let $\mathcal U$ be the multiplier-one subgroup in the
image of $\GTz$ on the free pro-$p$ group $P$ on $X,Y$.
Let $\mathcal U_n$ consist of pairs whose normalized conjugator belongs
to $\gamma_nP$.  Then
\[
 \mathcal U_2=\mathcal U_3=\mathcal U,
 \qquad \mathcal U_n/\mathcal U_{n+1}\simeq W_n(\Z_p),
 \qquad [\mathcal U_i,\mathcal U_j]\subset\mathcal U_{i+j}.
\]
Every element of the local multiplier-one group has a full gentle lift
of global multiplier one.
\end{lemma}
\begin{proof}
Write $Z=(XY)^{-1}$.  The ordinary lower-central graded Lie algebra of
$P$ is the free Lie algebra over $\Z_p$: integral Hall collection gives
the coordinates in each free nilpotent quotient, and every finite
$p$-group is nilpotent.  A leading normalized conjugator gives
$D_q(x)=0$, $D_q(y)=[y,q]$.  Preserving the three peripheral classes
and commuting with their permutations gives precisely the equations
defining $W_n$, as proved in Section~\ref{sec:dimensions}.  This also
proves injectivity of the leading map.

Conversely, start with an invariant outer special derivation of degree
$n$.  Normalize it to fix $x$, and lift its conjugators at $y,z$ to
$\gamma_nP$.  The prospective images of $X,Y,Z$ have product error in
$\gamma_{n+2}P$.  If the error lies in $\gamma_jP$, changing the two
conjugators in degree $j-1$ changes its degree-$j$ term by
\[
 (a,b)\longmapsto[a,y]+[b,-x-y].
\]
This map is onto: its image is $[L_{j-1},x]+[L_{j-1},y]=L_j$.
Successively cancel every error.  The corrections converge because
every finite $p$-group quotient kills sufficiently deep lower-central
terms.  The resulting endomorphism is the identity on the Frattini
quotient, hence is onto by the pro-$p$ basis theorem and is an
automorphism by finite generation.

The special outer automorphisms with trivial Frattini action form a
pro-$p$ group.  Their filtration is stable under the geometric $S_3$.
Fixed points lift through an onto equivariant map of pro-$p$ groups
when the acting finite group has order prime to $p$.  To see this,
the discrepancy of a lift is a nonabelian one-cocycle in the kernel;
the corresponding two complements in its semidirect product are
conjugate.  For a finite $p$-kernel, induction through a central
elementary abelian subgroup and averaging prove this assertion.
For a pro-$p$ kernel, compactness of the closed sets of possible
conjugating elements proves it from the finite cases.  Apply this to
the degree-$n$ quotient.  It changes only higher terms and gives an
$S_3$-invariant special outer automorphism with the required leading term.

Here is the explicit passage to the normalized group equations.
Let $\mathcal B=\widehat{C_2*C_3}$, with generators $s,r$, and set
$\tau=sr$, $F=\ker(\mathcal B\to S_3)$,
$X=\tau^{-2}$, $Y=sXs$.  Schreier rewriting gives $F=\widehat F_2$
on $X,Y$, and conjugation by $r$ cycles $X,Y,Z$.  The quotient
$B_p=\mathcal B/O^p(F)$ has kernel $P$ and outer quotient $S_3$.
The centre of $P$ is trivial, so
\[
 B_p\simeq\operatorname{Aut}(P)
           \mathbin{\times}_{\operatorname{Out}(P)} S_3.
\]
An invariant outer class therefore gives an automorphism of $B_p$
acting trivially on $S_3$.  Choose its representative to fix $X$.
Its image of $\tau$ has the form $X^a\tau$: both elements centralize
$X$, their quotient lies in $P$, and $C_P(X)=\langle X\rangle$.
Equality of their squares implies $2a=0$ in $\Z_p$, so it fixes $\tau$.
If $f_1=s\phi(s)$, then $f_1$ has abelianization $(a,-a)$.
Conjugating $\phi$ by $X^a$ replaces it by
$f=Y^af_1X^{-a}\in P'$.  The resulting formulas are
\[
 \phi(s)=sf,\qquad \phi(r)=f^{-1}r,
 \qquad X\mapsto X,\quad Y\mapsto f^{-1}Yf.
\]
The relations $s^2=r^3=1$ give, in $P$,
\[
 f(Y,X)f(X,Y)=1,\qquad f(Z,X)f(Y,Z)f(X,Y)=1.
\]
These are the multiplier-one gentle equations.  Centralizers in
$P/\gamma_{n+1}P$ show that the normalization changes no degree-$n$
outer term: an inner change fixing $x$ has zero homogeneous
conjugator in degree $n>1$.

Use this automorphism in the $p$-component of the maximal pronilpotent
pure quotient of $\mathcal B$, and the identity in every other component.
It preserves the positive $s$- and $r$-classes: their complements over
$C_2,C_3$ are conjugate by the same coprime argument.  It fixes $\tau$
exactly.  The full peripheral lifting theorem gives an automorphism
of $\mathcal B$ with these prescribed conjugators.  An inner correction
from the kernel makes it fix $\tau$ exactly.  Its normalized $f=s\phi(s)$
has zero abelianization at every prime, and the same torsion relations
give the full gentle equations.  Thus it has global multiplier one and
the prescribed local image.  This argument applies to an arbitrary
local solution, not only to a leading term.

Finally, composition of normalized conjugators gives the Ihara bracket
on the associated graded, hence the displayed commutator inclusion.
Normalization removes degree one and $W_2=0$ removes degree two.
\end{proof}

\begin{lemma}[Rational Lie algebra of a compact unipotent group]
\label{lem:compact-log}
Let $U$ be a closed subgroup of a unipotent matrix group over $\Q_p$.
Set $L=\operatorname{span}_{\Q_p}\log U$.  Then $L$ is a nilpotent Lie
algebra and $U$ contains $\exp(J)$ for some full $\Z_p$-lattice
$J\subset L$.  These statements apply to every closed subgroup of $U$.
\end{lemma}
\begin{proof}
All logarithms, exponentials and Baker--Campbell--Hausdorff expressions
are finite polynomials.  For $X=\log u$ and $Y=\log v$, the polynomial
$\operatorname{BCH}(aX,bY)$ belongs to $L$ for nonnegative integers
$a,b$.  Evaluating on a sufficiently large integer grid and inverting
the rational Vandermonde matrices shows that every coefficient belongs
to $L$.  Its $ab$-coefficient is $[X,Y]/2$, proving closure under brackets.

Choose a basis $X_i=\log u_i$ of $L$ from logarithms of elements of $U$.
For $t_i\in\Z_p$, continuity gives $u_i^{t_i}=\exp(t_iX_i)\in U$.
In these basis coordinates,
\[
 (t_i)\longmapsto\log(u_1^{t_1}\cdots u_d^{t_d})=t+R(t),
\]
where $R$ has no constant or linear term.  On a sufficiently small ball
$p^r\Z_p^d$, $R$ has Lipschitz constant at most $p^{-1}$ and maps that
ball into $p^{r+1}\Z_p^d$.  For each $y$ in the ball, contraction of
$t\mapsto y-R(t)$ solves $t+R(t)=y$.  Thus the exponential of the
corresponding full lattice is contained in $U$.  The proof uses only
closedness and applies verbatim to a closed subgroup.
\end{proof}

\begin{proof}[Proof of Theorem~\ref{thm:even-characters}]
Let $\mathcal G$ be the local pro-$p$ image of $\GTz$, and set
\[
 G=\mathcal G/\mathcal U_{m+1},\qquad
 U=\mathcal U/\mathcal U_{m+1}.
\]
The action on the integral Magnus algebra through augmentation degree
$m+1$ realizes $U$ faithfully by unipotent matrices.  Put
$L=\operatorname{span}_{\Q_p}\log U$.  Lemmas
\ref{lem:realization} and \ref{lem:compact-log} identify its degree
filtration with
\[
 L_n/L_{n+1}=W_n(\Q_p),\qquad 3\leq n\leq m.
\]
Indeed every leading group term gives the indicated Lie term, and the
open exponential lattice ensures that every filtered Lie vector, after
multiplication by a power of $p$, is a logarithm of a filtered group element.

Choose an element $g_0$ in the local Galois image whose multiplier is
$1+p$.  The action $A=\operatorname{Ad}(g_0)$ on $L_n/L_{n+1}$ is
$(1+p)^n$: conjugation of a degree-$n$ derivation by the scalar
action on degree one scales it by that power.  These eigenvalues are
distinct over $\Q_p$.  Applying the factors successively to the
filtration gives
\[
 \prod_{n=3}^{m}(A-(1+p)^n)=0.
\]
Thus $A$ is semisimple, and its eigenspaces give a graded splitting
\[
 L=\bigoplus_{n=3}^{m}L^{(n)},\qquad L^{(n)}\simeq W_n(\Q_p).
\]
Because $A$ is a Lie automorphism, brackets respect this grading and
are the Ihara brackets on the leading terms.

Choose the full space of degree-$m$ linear functionals on
$L/[L,L]$, extending each by zero on the other eigenspaces.  They give
a map $\ell:L\to\Q_p^{h_m}$ annihilating $[L,L]$.
Consequently $u\mapsto\ell(\log u)$ is an additive homomorphism.
Its image is a compact $\Z_p$-submodule spanning $\Q_p^{h_m}$, and
hence is a full free lattice.  Choosing a basis of this image gives a
continuous surjection
\[
 \Phi:U\longrightarrow\Z_p^{h_m}.
\]
Conjugation on $L/[L,L]$ factors through
$G/U=\Z_p^\times$, so every such conjugation commutes with $A$ on
that abelianization.  It preserves its weight spaces.  On the top
weight it acts as on the leading quotient, by the scalar
$\lambda_p(g)^m$.  Therefore
\begin{equation}
 \Phi(gug^{-1})=\lambda_p(g)^m\Phi(u).
 \label{eq:equivariance}
\end{equation}

Let $P_0$ be the image of $\GT$ in $G$ and put $P_1=P_0\cap U$.
By Lemma~\ref{lem:compact-log},
$L_P=\operatorname{span}_{\Q_p}\log P_1$ is a Lie subalgebra and
$P_1$ contains an exponential neighborhood in $L_P$.
Since $g_0$ is Galois, it belongs to $P_0$; thus $L_P$ is
$A$-stable.  The polynomial projector $e_m$ onto $L^{(m)}$ preserves
$L_P$.  For $h\in P_1$, put $X=e_m\log h$.  Then
$X\in L_P\cap L_m$, and for some $r$, $\exp(p^rX)\in P_1$.
This is the image of a full pentagon element whose local conjugator
has no terms below weight $m$.  Its leading term therefore satisfies
the linear pentagon.  By \eqref{eq:even-pentagon}, it is a bracket in
$W_m(\Q_p)$, so $\ell(p^rX)=0$ and $\ell(X)=0$.
This proves $\Phi(P_1)=0$.

Reduce $\Phi$ modulo $p^k$.  The divisibility hypothesis gives
$a^m=1$ in $(\Z/p^k\Z)^\times$ for every $a\in\Z_p^\times$,
so \eqref{eq:equivariance} makes this reduction, denoted $\theta$,
invariant under all of $G$.  The multiplier on $P_0$ is onto,
because it contains the Galois image.  Hence $G=UP_0$.
Define
\[
 \chi(up)=\theta(u)\qquad(u\in U,\ p\in P_0).
\]
Two such factorizations differ by an element of $P_1$, so the value
is well defined.  Invariance of $\theta$ proves the homomorphism law.
The multiplication map from the compact space $U\times P_0$ onto
$G$ is a quotient map, which proves continuity.  The homomorphism
vanishes on $P_0$ and is onto.  Composing with the full gentle action
proves the asserted quotient.  The final clause of
Lemma~\ref{lem:realization} supplies all values on global multiplier one.
\end{proof}

\begin{corollary}[All finite odd-primary orders]\label{cor:all-orders}
Put $Q=\GTz/\overline{\langle\!\langle\GT\rangle\!\rangle}$.
There is a continuous surjection
\[
 Q^{\mathrm{ab}}\longrightarrow
       \prod_{\substack{p\geq5\\p\ \mathrm{prime}}}
          \prod_{k,j\geq1}\Z/p^k\Z.
\]
For every prime $p\geq5$, the vector space
$\Hom_{\cts}(Q,\F_p)$ has dimension exactly $\aleph_0$.
\end{corollary}
\begin{proof}
We first prove that $h_m$ tends to infinity along every even arithmetic
progression.  Let $a_n$ and $b_n$ be as in
\eqref{eq:dimension}--\eqref{eq:raw}, and put $r_n=a_n2^{-n}$.
The identity contribution to the Witt formula gives
\[
 a_n=\frac{2^n}{3n(n+1)}+O(2^{n/2}),\qquad
 r_n\sim\frac1{3n^2}.
\]
All nonidentity contributions and all nontrivial M\"obius divisors
have exponential size at most the stated error.
Write $l_n=\dim L_n$.  The dimension of
$\operatorname{OutSDer}_n$ is $2l_n-l_{n+1}$, whence
\[
 \sum_{n=2}^N r_n\leq
 \sum_{n=2}^N\frac{2l_n-l_{n+1}}{2^n}
       =\frac12-\frac{l_{N+1}}{2^N}.
\]
At $n=5$, the upper-bound dimension is three whereas $a_5=1$.
Consequently $S=\sum_{n\geq2}r_n\leq7/16$.
For any fixed $M$, the terms with $i\leq M$ or $n-i\leq M$ in
$\sum_i r_ir_{n-i}/r_n$ tend to $2\sum_{i\leq M}r_i$.
In the remaining half-interval, $r_{n-i}/r_n$ is bounded uniformly
and $\sum_{i>M}r_i=O(M^{-1})$.  Letting $M$ increase proves
\[
 \frac{b_n}{a_n}\longrightarrow S<1,
 \qquad h_n\geq a_n-b_n\sim(1-S)\frac{2^n}{3n^2}.
\]
The diagonal subtraction in the unordered pair count is exponentially
smaller and does not affect the limit.

Fix $p,k$.  Theorem~\ref{thm:even-characters}, with arbitrarily large
multiples of $(p-1)p^{k-1}$, gives arbitrarily many characters to
$\Z/p^k\Z$ whose reductions modulo $p$ are independent already on
global multiplier one.  Thus the space of these reductions is
infinite-dimensional.  Enumerate all pairs $(k,j)$.  At each stage
choose a character whose modulo-$p$ restriction lies outside the
finite span of the preceding restrictions.  For any finite collection,
let $H$ be its image in the target product $T$ of cyclic $p$-groups.
Independence gives $H+pT=T$.  Hence the finite group $T/H$ equals
$p(T/H)$ and is zero.  The finite product map is onto.

For different primes the primary targets have coprime orders, so a
subdirect product is their product.  Every finite projection of the
displayed infinite map is therefore onto.  Compactness proves
surjectivity onto the whole product.  Finally, the gentle group is
countably based: a finitely generated free profinite group has a
countable cofinal family of characteristic open subgroups.  Its
automorphism group, and the closed gentle subgroup, inherit this
property.  There are only countably many clopen subsets of a compact
space with a countable clopen basis, hence only countably many maps
to $\F_p$.  The lower bound already proved gives equality.
\end{proof}

\begin{remark}
The corollary specifies a quotient of the abelianization.  It does not
identify the whole abelianization, prove that the original quotient is
abelian, or construct a compatible $\Z_p$ quotient.  All its characters
vanish on the full pentagon group, so they give no separation of that
group from the absolute Galois group.
\end{remark}
